\documentclass[12pt,a4paper]{article}
\usepackage[utf8]{inputenc}
\usepackage[italian]{babel}
\usepackage{geometry}
\usepackage{graphicx}
\usepackage{listings}
\usepackage{xcolor}
\usepackage{hyperref}
\usepackage{booktabs}
\usepackage{tabularx}
\usepackage{fancyhdr}
\usepackage{float}
\usepackage{amsmath}

\geometry{margin=2.5cm}

\hypersetup{
    colorlinks=true,
    linkcolor=blue!60!black,
    urlcolor=blue!60!black,
    citecolor=blue!60!black
}

\lstset{
    basicstyle=\ttfamily\small,
    breaklines=true,
    frame=single,
    backgroundcolor=\color{gray!5},
    showstringspaces=false,
    tabsize=2
}

\pagestyle{fancy}
\fancyhf{}
\fancyhead[L]{\small Shipping on the Air with Patterns}
\fancyhead[R]{\small SAP - Assignment 2}
\fancyfoot[C]{\thepage}

\title{
    \textbf{Shipping on the Air with Patterns} \\[0.5em]
    \large Assignment \#02 - Software Architecture and Platforms \\[0.3em]
    \normalsize A.Y. 2025/2026
}
\author{Fabio Vincenzi}
\date{\today}

\begin{document}

\maketitle
\tableofcontents
\newpage

% ============================================================
\section{Introduzione}

Questo documento descrive, a livello architetturale, come il sistema \emph{Shipping on the Air}
progettato nell'Assignment~1 sia stato raffinato applicando un sottoinsieme dei microservices
patterns discussi nel corso.

Il dominio e le funzionalit\`a non cambiano: il sistema resta composto da \texttt{order-service},
\texttt{delivery-service} e \texttt{drone-service}, ai quali si aggiunge l'\texttt{api-gateway}.
Ci\`o che cambia \`e il modo in cui i servizi comunicano con l'esterno, vengono osservati a runtime,
persistono il proprio stato e vengono verificati.

Ogni pattern \`e presentato indicando il problema che risolve nel contesto specifico del progetto,
la soluzione adottata e le conseguenze, incluse quelle negative. Dove possibile le affermazioni
sono accompagnate da misure raccolte sul sistema in esecuzione.

\paragraph{Pattern applicati.} Oltre ai tre richiesti dalla consegna --- API Gateway,
Observability Patterns (Health Check API e Application Metrics) ed Event Sourcing applicato a un
servizio --- sono stati scelti \textbf{CQRS} e \textbf{Service Discovery}. La scelta non \`e
arbitraria: come argomentato nella sezione~\ref{sec:cqrs}, CQRS non \`e un'aggiunta facoltativa ma
una conseguenza tecnica diretta dell'Event Sourcing.

% ============================================================
\section{Quality Attribute Scenarios}
\label{sec:qas}

Gli attributi di qualit\`a individuati nell'Assignment~1 sono stati riformulati in due scenari
concreti, secondo la struttura in sei parti vista a lezione. Entrambi riguardano la stessa
propriet\`a architetturale --- la \emph{degradazione controllata} --- ma sotto due stimoli diversi:
il guasto di un servizio e l'esaurimento di una risorsa.

\begin{lstlisting}[caption={QAS-A --- Availability: guasto di un servizio a valle}]
Source            un guasto infrastrutturale
Stimulus          il drone-service diventa irraggiungibile
Environment       operazioni normali, sistema deployato in container
Artifact          delivery-service e piattaforma di deployment
Response          il guasto viene rilevato e il sistema continua ad
                  accettare ordini e a programmare consegne, senza
                  assegnare droni
Response Measure  rilevamento entro 15 s; nessun fallimento propagato
                  al cliente
\end{lstlisting}

\begin{lstlisting}[caption={QAS-B --- Saturazione: flotta di droni esaurita}]
Source            domanda superiore alla capacita del sistema
Stimulus          arriva un ordine quando tutti i droni sono occupati
Environment       operazioni normali, flotta satura
Artifact          delivery-service e drone-service
Response          l'ordine viene accettato e la consegna programmata,
                  in attesa che un drone si liberi
Response Measure  nessun ordine rifiutato; saturazione osservabile
                  come drones_available = 0
\end{lstlisting}

La verifica di questi scenari \`e discussa nella sezione~\ref{sec:osservabilita-qas}, dopo aver
introdotto gli observability pattern che li rendono misurabili.

% ============================================================
\section{Deployment}

Il sistema adotta il pattern \textbf{Service per Container}: ogni microservizio viene impacchettato
in un'immagine Docker separata e il sistema completo \`e descritto da un file Docker Compose.

Ogni servizio ha un \texttt{Dockerfile} \textbf{multi-stage}: un primo stage basato su JDK e Gradle
esegue \texttt{installDist}, un secondo stage contiene solo il JRE e la distribuzione prodotta.
L'immagine finale non contiene quindi n\'e il JDK n\'e il build tool. Il contesto di build \`e la
radice del monorepo, perch\'e tutti i servizi dipendono dal modulo \texttt{common}.

Il file \texttt{docker-compose.yml} definisce la rete logica \texttt{shipping\_network}, alla quale
tutti i servizi sono collegati, e imposta per ciascuno la policy \texttt{restart: unless-stopped} e
un \texttt{healthcheck} (sezione~\ref{sec:healthcheck}). La configurazione dipendente dall'ambiente
\`e iniettata tramite variabili d'ambiente, secondo il pattern \textbf{Externalized Configuration}:
i launcher leggono host e porta dei servizi con cui devono comunicare, usando \texttt{localhost}
come valore di fallback. Lo stesso artefatto funziona quindi sia in esecuzione locale sia in
container, senza modifiche al codice.

% ============================================================
\section{Patterns}

\subsection{API Gateway}

\paragraph{Problema.} Senza un punto di ingresso unico, un client dovrebbe conoscere indirizzo e
numero dei servizi, effettuare pi\`u chiamate per popolare una singola schermata e ogni servizio
dovrebbe reimplementare le funzionalit\`a trasversali.

\paragraph{Soluzione.} \`E stato aggiunto il modulo \texttt{api-gateway} (porta 8080, prefisso
\texttt{/api/v1}), unico componente esposto all'esterno. Il \texttt{GatewayController} instrada le
richieste verso il servizio competente attraverso tre proxy (\texttt{OrderServiceProxy},
\texttt{DeliveryServiceProxy}, \texttt{DroneServiceProxy}), definiti come porte in uscita
nell'application layer e implementati nell'infrastructure layer.

\paragraph{API Composition.} L'endpoint di tracking \texttt{GET /api/v1/tracking/:orderId} non si
limita a instradare: interroga l'\texttt{order-service} e il \texttt{delivery-service} e
\textbf{unisce le due risposte} in un unico oggetto \texttt{TrackingInfo}. \`E un'applicazione del
pattern \emph{API Composition}, necessario perch\'e con un database per servizio non \`e possibile
eseguire join tra dati di servizi diversi. Il client effettua una chiamata sola invece di due.

Poich\'e i proxy sono sincroni e bloccanti, le chiamate a valle sono eseguite tramite
\texttt{executeBlocking} con \texttt{ordered = false}, in modo da non bloccare l'event loop di
Vert.x.

\paragraph{Conseguenze.} Il gateway semplifica i client e disaccoppia il loro codice dalla struttura
interna del sistema, ma introduce un componente che deve essere reso altamente disponibile e che
va aggiornato ogni volta che cambia un'API a valle.

\subsection{Observability Patterns}

\subsubsection{Health Check API}
\label{sec:healthcheck}

Tutti e quattro i componenti espongono un endpoint \texttt{/health} che restituisce lo stato del
servizio nel formato MicroProfile Health (\texttt{status} e lista di \texttt{checks}).

L'endpoint non serve solo alla diagnosi manuale: \`e consumato dall'attributo \texttt{healthcheck}
del Docker Compose file, che lo interroga periodicamente e marca il container come
\texttt{unhealthy} in caso di fallimento. Combinato con la policy di restart, questo permette alla
piattaforma di reagire autonomamente a un servizio bloccato.

\subsubsection{Application Metrics}

Ogni componente espone metriche in formato Prometheus su una \textbf{porta dedicata}, distinta da
quella di servizio, tramite l'\texttt{HTTPServer} della libreria \texttt{io.prometheus}. Un
container \texttt{prometheus} raccoglie i quattro target con \texttt{scrape\_interval: 10s}.

\begin{table}[H]
\centering
\begin{tabularx}{\textwidth}{llX}
\toprule
\textbf{Componente} & \textbf{Metrica} & \textbf{Tipo e significato} \\
\midrule
order-service & \texttt{orders\_placed\_total} & counter: ordini creati \\
delivery-service & \texttt{deliveries\_in\_progress} & gauge: consegne in corso \\
drone-service & \texttt{drones\_available} & gauge: droni liberi \\
api-gateway & \texttt{gateway\_requests\_total} & counter: richieste ricevute \\
\bottomrule
\end{tabularx}
\end{table}

La scelta del tipo di metrica riflette la natura del fenomeno osservato: un \emph{counter} per
eventi cumulativi che possono solo crescere, un \emph{gauge} per quantit\`a che salgono e scendono.
Sui counter va sempre usato \texttt{rate()}, poich\'e il valore assoluto si azzera al riavvio del
processo.

\paragraph{Come sono raccolte.} Per i tre servizi di dominio la strumentazione non \`e inserita nei
controller ma realizzata tramite \textbf{observer}: l'application layer definisce una porta in
uscita (ad esempio \texttt{DeliveryServiceObserver}) e l'infrastructure layer ne fornisce un adapter
(\texttt{PrometheusDeliveryServiceObserver}). Il dominio non conosce quindi Prometheus e resta
osservabile senza esserne dipendente. Il gateway costituisce un'eccezione consapevole: non avendo un
core di dominio, la sua metrica \`e puramente infrastrutturale e viene incrementata da un handler di
routing.

\subsection{Event Sourcing}
\label{sec:eventsourcing}

\paragraph{Scelta del servizio.} L'Event Sourcing \`e stato applicato al \texttt{delivery-service}.
\`E un pattern invasivo, che richiede di ristrutturare la business logic, quindi va applicato dove
produce valore: la consegna ha un ciclo di vita a stati articolato
(\texttt{SCHEDULED} $\rightarrow$ \texttt{DRONE\_ASSIGNED} $\rightarrow$ \texttt{IN\_TRANSIT}
$\rightarrow$ \texttt{DELIVERED}/\texttt{FAILED}) e una storia che ha valore di per s\'e --- dove si
trovava il drone in un dato istante, quando \`e stato assegnato. Applicarlo anche agli altri servizi
avrebbe aggiunto complessit\`a senza beneficio.

\paragraph{Ristrutturazione dell'aggregate.} Ogni metodo di comando \`e stato separato in due:

\begin{itemize}
  \item i metodi \texttt{processXxx()} validano la richiesta, \textbf{non modificano lo stato} e
        restituiscono la lista di eventi da applicare; sono gli unici che possono fallire;
  \item i metodi \texttt{apply()} aggiornano lo stato a partire da un evento e non possono fallire,
        perch\'e un evento descrive un fatto gi\`a accaduto.
\end{itemize}

Tutte e cinque le transizioni emettono ora un evento (\texttt{DeliveryScheduled},
\texttt{DroneAssigned}, \texttt{TransitStarted}, \texttt{DronePositionUpdated},
\texttt{DeliveryCompleted}, \texttt{DeliveryFailed}); in precedenza solo due lo facevano.

Una conseguenza non ovvia riguarda il \textbf{determinismo}: poich\'e lo stato viene ricostruito
riapplicando gli eventi, i metodi \texttt{apply()} non possono contenere sorgenti di
non-determinismo. L'istante di creazione, prima calcolato con \texttt{Instant.now()} nel costruttore,
\`e stato spostato \textbf{dentro} l'evento \texttt{DeliveryScheduled}, in modo che ogni replay
riproduca esattamente lo stesso stato.

\paragraph{Event store.} La porta \texttt{DeliveryEventStore} espone \texttt{append}, \texttt{load} e
\texttt{loadFrom}: non esiste alcuna operazione di modifica o cancellazione, quindi l'append-only
\`e garantito dalla forma dell'interfaccia e non dalla disciplina di chi la usa.

Il metodo \texttt{append} riceve una \texttt{expectedVersion}, cio\`e la versione a cui il chiamante
ha letto l'aggregate, e rifiuta la scrittura se lo stream \`e nel frattempo avanzato. \`E un
controllo di \textbf{concorrenza ottimistica} che impedisce a due comandi concorrenti di credersi
entrambi l'ultimo, e costituisce la differenza sostanziale tra un event store e una semplice coda di
messaggi.

\paragraph{Snapshot.} Poich\'e il \texttt{delivery-service} riceve aggiornamenti di posizione ad alta
frequenza, lo stream di una consegna cresce indefinitamente e il costo di ogni lettura crescerebbe
con esso. Viene quindi persistito uno snapshot dello stato ogni 50 eventi: il caricamento parte dallo
snapshot pi\`u recente e riapplica soltanto gli eventi successivi.

Misurato in locale su uno stream di 1003 eventi, con media su 2000 caricamenti dopo warm-up:

\begin{table}[H]
\centering
\begin{tabular}{lc}
\toprule
& \textbf{Tempo per \texttt{findById}} \\
\midrule
senza snapshot (replay di 1003 eventi) & 0,013 ms \\
con snapshot (replay dei soli eventi successivi) & 0,004 ms \\
\bottomrule
\end{tabular}
\end{table}

Il guadagno misurato \`e di circa 3,6$\times$. Il rapporto \`e contenuto perch\'e event store e
snapshot sono in memoria e gli eventi hanno un payload minimo, quindi si sta misurando il solo costo
del \emph{fold}; con una persistenza su database il divario crescerebbe, dovendo leggere 1003 righe
invece di uno snapshot e pochi eventi. Il punto architetturale, comunque, non \`e il fattore ma il
fatto che il costo passi da $O(\text{numero di eventi})$ a $O(\text{eventi dopo lo snapshot})$, cio\`e
da crescente a costante.

Lo snapshot \`e \textbf{stato derivato}: eliminarlo lascia il sistema corretto, solo pi\`u lento.
La propriet\`a \`e verificata dal test \texttt{droppingEverySnapshotKeepsTheStateCorrect}, che carica
la stessa consegna con e senza snapshot e confronta i due stati risultanti.

\subsection{CQRS}
\label{sec:cqrs}

\paragraph{Perch\'e \`e necessario.} Un event store risponde soltanto per chiave primaria, cio\`e per
identificatore di consegna. Il \texttt{DeliveryRepository}, per\`o, espone anche
\texttt{findByOrderId} e \texttt{findByDroneId}: due interrogazioni su campi che chiave non sono.
Dopo l'introduzione dell'Event Sourcing queste operazioni non sarebbero pi\`u implementabili se non
ricostruendo l'intero insieme delle consegne a ogni richiesta.

CQRS non \`e quindi un pattern aggiunto per scelta estetica, ma la conseguenza diretta della
decisione presa nella sezione~\ref{sec:eventsourcing}: \textbf{\`e l'Event Sourcing a imporlo}.

\paragraph{Soluzione.} I due lati sono separati anche strutturalmente:

\begin{itemize}
  \item il \textbf{write model} \`e costituito dall'aggregate e dall'event store, append-only e
        fonte di verit\`a;
  \item il \textbf{read model} \`e la porta \texttt{DeliveryLookupView} con il relativo adapter, che
        mantiene due indici \texttt{orderId $\rightarrow$ deliveryId} e
        \texttt{droneId $\rightarrow$ deliveryId};
  \item il \texttt{DeliveryProjector} traduce gli eventi in aggiornamenti della vista ed \`e
        \textbf{l'unico componente autorizzato a scriverci}.
\end{itemize}

Il repository scrive attraverso il projector e legge attraverso la vista: i comandi non toccano mai
il read model e le interrogazioni non attraversano mai il modello di dominio.

La vista contiene deliberatamente i soli identificatori e non lo stato della consegna. In questo modo
l'unica informazione duplicata \`e una corrispondenza che non cambia mai nel tempo: nel peggiore dei
casi una consegna non viene trovata, mai vengono restituiti dati errati.

\paragraph{Ricostruibilit\`a.} Il metodo \texttt{rebuildLookupView()} svuota la vista e la ricostruisce
riapplicando tutti gli eventi memorizzati, usando la stessa funzione di proiezione impiegata a regime.
Anche il read model, come lo snapshot, \`e dato derivato: il test
\texttt{theViewCanBeDroppedAndRebuiltFromTheEvents} lo dimostra.

\paragraph{Consistenza.} Nell'implementazione realizzata la proiezione \`e invocata in modo sincrono
subito dopo la scrittura degli eventi, quindi la vista \`e sempre allineata e non si manifesta il
problema del \emph{read-your-own-writes}. In cambio non si ottiene il disaccoppiamento completo: il
passaggio a una proiezione asincrona richiederebbe di spostare la chiamata al projector dietro una
sottoscrizione all'event store, \textbf{senza modificare n\'e la vista n\'e il repository}, proprio
perch\'e la vista \`e una porta separata.

\subsection{Service Discovery}

\paragraph{Problema.} In un deployment a container gli indirizzi dei servizi sono dinamici, quindi non
possono essere scritti nel codice n\'e in un file di configurazione statico.

\paragraph{Soluzione.} Il pattern \`e realizzato \textbf{a livello di deployment}: Docker fornisce
nativamente un service registry e un DNS resolver interno, assegnando a ogni servizio un nome sulla
rete \texttt{shipping\_network}. Si tratta del modello \emph{server-side discovery}: il chiamante usa
sempre lo stesso nome logico e l'infrastruttura risolve verso un'istanza attiva.

Un dettaglio determinante \`e che le variabili d'ambiente iniettate dal Compose file puntano ai
\textbf{nomi di servizio} (\texttt{DELIVERY\_HOST=delivery-service}) e non ai
\texttt{container\_name} (\texttt{delivery-01}). \`E la differenza tra discovery effettiva e indirizzo
fisso: il nome di servizio viene risolto dinamicamente e, in presenza di pi\`u repliche, distribuito
tra tutte; il nome del container punterebbe sempre alla stessa istanza.

Il pattern si integra con l'Externalized Configuration: lo stesso artefatto usa \texttt{localhost} in
esecuzione locale e il nome logico in container. I limiti attuali di questa configurazione sono
discussi nella sezione~\ref{sec:limiti}.

% ============================================================
\section{Testing}

La strategia di testing segue la test pyramid, con almeno un esempio per ciascun livello. I livelli
sono distinti nella struttura dei package, in modo che la piramide sia leggibile dall'organizzazione
del codice e non solo da questa descrizione.

\subsection{Unit}

Gli unit test sono organizzati per strato dell'architettura esagonale, seguendo la distinzione tra
test \emph{solitary} (dipendenze sostituite da mock) e \emph{sociable} (dipendenze reali).

\begin{itemize}
  \item \textbf{application}: \texttt{DeliveryServiceTest} verifica che il completamento di una
        consegna notifichi l'order-service e rilasci il drone; le porte verso gli altri servizi sono
        sostituite da stub. \`E un test \emph{solitary}.
  \item \textbf{infrastructure}: \texttt{DeliveryEventStoreTest}, \texttt{DeliverySnapshotTest} e
        \texttt{DeliveryProjectionTest} verificano gli adapter di persistenza --- append-only,
        ordinamento, rifiuto delle scritture concorrenti, replay, snapshot e ricostruzione della vista
        --- usando gli oggetti di dominio reali. Sono test \emph{sociable}.
  \item \textbf{controller}: \texttt{DeliveryControllerTest} verifica rotte, parsing del payload,
        codici di stato e formato JSON, con un \texttt{DeliveryServiceMock} al posto del servizio.
        \`E un test \emph{solitary}: il soggetto \`e una sola classe, e la porta HTTP \`e soltanto il
        modo in cui un adapter di questo tipo pu\`o essere invocato.
\end{itemize}

Un test merita una nota particolare: \texttt{DeliveryEventStoreTest} verifica che l'istante di
creazione sopravviva identico a un replay, proteggendo cos\`i la propriet\`a di determinismo
descritta nella sezione~\ref{sec:eventsourcing}.

\subsection{Integration}

\texttt{DroneServiceProxyTest} verifica il contratto tra due servizi: il \texttt{DroneController}
reale del drone-service viene avviato su una porta di test con un servizio applicativo reale su
repository in memoria, e vi si punta contro il \texttt{DroneServiceProxy} reale del delivery-service.
Ci\`o che viene verificato non \`e la logica di business, ma che URL, query parameter, formato JSON e
codici di stato dei due lati coincidano.

Essendo il progetto un monorepo Gradle, il test dichiara semplicemente
\texttt{testImplementation project(':drone-service')}: non \`e stato necessario duplicare classi tra i
moduli, e la dipendenza esiste soltanto in ambito di test.

Un caso di questo gruppo verifica una propriet\`a rilevante per la sezione~\ref{sec:osservabilita-qas}:
con il servizio a valle irraggiungibile, il proxy restituisce un risultato vuoto anzich\'e propagare
l'eccezione.

\subsection{Component}

Gli acceptance test del singolo servizio sono scritti in Gherkin ed eseguiti con Cucumber. Il
delivery-service viene avviato \textbf{per intero} --- controller, application layer e persistenza a
eventi --- su una porta di test, e gli step lo interrogano \textbf{attraverso la sua API HTTP}, come
farebbe un client. Solo le porte verso gli altri servizi sono sostituite da stub.

Il confine del test \`e quindi l'API pubblica del servizio, non le sue classi: gli step manipolano
oggetti JSON e non hanno accesso al modello di dominio. Ne consegue che una modifica interna non
rompe questi test, mentre una modifica al contratto pubblico s\`i --- che \`e esattamente il
comportamento desiderato per un acceptance test.

Si tratta di component test \emph{in-process}: la rete e il protocollo HTTP sono reali, mentre la
verifica del packaging in container \`e demandata al livello end-to-end.

\subsection{End-to-end}

Il modulo separato \texttt{system-end-to-end} contiene uno user journey che attraversa l'intero
sistema passando esclusivamente dall'API gateway: registrazione di un drone e di un utente, creazione
di un ordine, conferma, e verifica che il tracking riporti una consegna con un drone assegnato. Il
percorso esercita tutti e quattro i componenti e l'API Composition del gateway.

Poich\'e richiede il sistema in esecuzione, questo modulo \`e escluso dalla build ordinaria ed \`e
attivato da una propriet\`a esplicita:

\begin{lstlisting}[caption={Esecuzione dei test end-to-end}]
docker compose up -d
./gradlew :system-end-to-end:test -Pe2e
\end{lstlisting}

\subsection{Fitness function}

Oltre ai quattro livelli, ogni servizio include un \texttt{ArchitectureTest} basato su ArchUnit, che
verifica che il dominio non dipenda da application e infrastructure, che le porte risiedano
nell'application layer e che gli adapter risiedano nell'infrastructure layer. Non appartiene alla test
pyramid: \`e una \emph{fitness function}, cio\`e un controllo automatico di una caratteristica
architetturale, e come tale uno strumento di governance dell'architettura.

% ============================================================
\section{Observability al servizio dei Quality Attribute Scenarios}
\label{sec:osservabilita-qas}

Un Quality Attribute Scenario \`e verificabile solo se la sua \emph{Response Measure} \`e osservabile.
Senza strumentazione uno scenario resta un'intenzione: \`e questo il ruolo degli observability
pattern rispetto ai requisiti di qualit\`a. La relazione \`e analoga a quella delle fitness function,
con la differenza che queste verificano caratteristiche architetturali a tempo di build, mentre
l'observability le verifica a runtime.

Entrambi gli scenari della sezione~\ref{sec:qas} sono stati tradotti in test eseguibili
(\texttt{QualityAttributeScenariosTest}) che provocano la condizione avversa e leggono le metriche.

\subsection{QAS-A --- Availability}

Gli observability pattern coinvolti sono due. L'\textbf{Health Check API} rende interrogabile lo stato
di ciascun servizio; le \textbf{Application Metrics} lo trasformano in una serie temporale, perch\'e
Prometheus genera automaticamente per ogni target la metrica \texttt{up}, pari a 1 finch\'e lo scrape
riesce e 0 appena fallisce. \`E il punto in cui i due pattern si saldano.

Il test arresta il container del drone-service, attende che la metrica cambi e verifica il
comportamento del sistema in quello stato:

\begin{lstlisting}[caption={Esito osservato}]
up{job="monitoring-drone-service"} = 0     (gli altri tre restano a 1)
guasto rilevato in 8,157 s

ordine creato -> PENDING
confirm       -> HTTP 200
tracking      -> hasDelivery: true, deliveryStatus: "SCHEDULED",
                 droneId: null
\end{lstlisting}

La consegna viene programmata comunque, semplicemente priva di drone: il guasto non si propaga al
cliente. Il tempo di rilevamento misurato \`e coerente con lo \texttt{scrape\_interval} di 10~s: la
latenza di rilevamento \`e determinata dalla frequenza di campionamento, non dal servizio. La
\emph{Response Measure} dichiarata (entro 15~s) \`e quindi soddisfatta, ma \`e importante osservare
che essa dipende da una scelta di configurazione dell'infrastruttura di monitoraggio.

\`E utile notare come questa propriet\`a sia coperta due volte, in due modi complementari: il test di
integrazione \texttt{anUnreachableServiceYieldsNoDrone} la \textbf{garantisce} prima del deploy,
verificando che il proxy restituisca un risultato vuoto; la metrica la \textbf{osserva} in esercizio.

\subsection{QAS-B --- Saturazione}

Qui l'observability pattern rilevante \`e la sola Application Metrics, attraverso il gauge
\texttt{drones\_available}. La scelta del tipo di metrica \`e essa stessa parte della soluzione: un
gauge pu\`o salire e scendere ed \`e quindi adatto a rappresentare una risorsa esauribile, mentre un
counter potrebbe solo indicare quanti droni siano stati registrati complessivamente.

Il test crea e conferma ordini finch\'e il gauge non raggiunge lo zero, quindi effettua un ulteriore
ordine e verifica che venga comunque accettato e programmato, in attesa di un drone. La forma della
risposta \`e la stessa del QAS-A --- \texttt{SCHEDULED} con \texttt{droneId} nullo --- ma la causa \`e
l'esaurimento di una risorsa anzich\'e un guasto.

% ============================================================
\section{Limiti noti}
\label{sec:limiti}

\paragraph{I servizi non sono replicabili orizzontalmente.} La configurazione attuale prevede
un'istanza per servizio. Tentando di avviare pi\`u repliche emergono tre ostacoli, di gravit\`a
crescente: i \texttt{container\_name} espliciti impediscono lo scaling, perch\'e i nomi devono essere
univoci; la pubblicazione di porte fisse sull'host causa un conflitto di binding; e soprattutto i
servizi mantengono lo stato in repository in memoria, per cui repliche diverse avrebbero stati
divergenti. I primi due ostacoli sono scelte di leggibilit\`a in fase di sviluppo, facilmente
reversibili; il terzo \`e strutturale. La strada per lo scaling passa dall'estrarre lo stato in uno
store condiviso, operazione che nel delivery-service richiederebbe di sostituire il solo adapter, dato
che l'event store \`e gi\`a definito dietro una porta.

\paragraph{La persistenza \`e in memoria.} L'event store, gli snapshot e la vista di lettura sono
implementati con strutture dati in memoria, coerentemente con gli esempi visti a lezione. Ne consegue
che la storia degli eventi non sopravvive al riavvio di un container, il che \`e in tensione con uno
dei benefici dell'Event Sourcing, l'audit trail permanente. Anche in questo caso il costo del
cambiamento \`e limitato alla sostituzione dell'adapter.

\paragraph{Non tutti i QAS sono misurabili.} Lo scenario di performance formulato nell'Assignment~1
--- tracking con latenza inferiore al secondo --- non \`e verificabile con la strumentazione attuale,
composta da counter e gauge, che indicano quante richieste siano state servite ma non quanto tempo
abbiano impiegato. Servirebbe un istogramma delle durate sul gateway, dal quale ricavare i percentili.
La conseguenza generale merita di essere esplicitata: \textbf{lo strumento di misura determina quali
attributi di qualit\`a \`e possibile promettere}.

\paragraph{Event Sourcing parziale.} L'event store realizzato copre il lato database della
definizione, ossia l'inserimento e il recupero degli eventi per chiave, ma non espone un meccanismo di
sottoscrizione: le notifiche verso gli altri servizi restano effettuate esplicitamente dall'application
layer. Non sono inoltre affrontate l'evoluzione dello schema degli eventi n\'e la cancellazione dei
dati personali, entrambe rilevanti in un sistema destinato alla produzione.

% ============================================================
\section{Conclusioni}

L'applicazione dei pattern ha modificato il sistema lungo tre direzioni distinte.

Verso l'esterno, l'API Gateway ha reso il sistema utilizzabile attraverso un unico punto di ingresso,
nascondendo ai client la decomposizione interna; l'API Composition ha reso possibile una query che
attraversa due servizi, altrimenti impraticabile con un database per servizio.

A runtime, gli observability pattern hanno trasformato requisiti di qualit\`a in propriet\`a
misurabili. \`E l'aspetto pi\`u significativo del lavoro: senza strumentazione, uno scenario come
``il guasto di un servizio non deve propagarsi'' non sarebbe stato n\'e verificabile n\'e
dimostrabile, mentre \`e stato possibile misurarne il rilevamento in 8,157~s e osservare il sistema
continuare a operare in modo degradato.

Sulla persistenza, l'Event Sourcing ha sostituito la memorizzazione dello stato con quella della sua
storia, ottenendo un registro immutabile e la possibilit\`a di ricostruire lo stato a qualunque
istante, al prezzo di una ristrutturazione della business logic e della necessit\`a di introdurre
CQRS per le interrogazioni non basate su chiave. Quest'ultimo punto \`e forse l'insegnamento pi\`u
utile dell'intero assignment: i pattern non sono elementi indipendenti da collezionare, ma soluzioni
che si richiamano l'una con l'altra, e adottarne uno significa spesso accettarne un secondo.

Infine, la test pyramid ha permesso di distribuire la verifica su livelli con obiettivi diversi ---
dalla singola classe fino allo user journey attraverso il sistema completo --- e di collegare i test
alle metriche: le stesse propriet\`a garantite dai test prima del deploy sono osservate dalle metriche
durante l'esercizio.

\end{document}
